\documentclass[11pt]{article}

% ---------------------------------------------------------------------------
%  L9 — The cut-count cancellation identity.
%  Category L (entirely machine-written); see paper/README.md and
%  docs/publication-split.md.
%
%  NOVELTY STATUS: UNCHECKED.  No literature-priority pass has been run.  The
%  identity is of a well-populated type (transfer-matrix connectivity counting
%  by colourings), so a collision is more likely here than in most of this
%  series, and nothing is claimed to be new.
%
%  Source material, all in this repository:
%    docs/proofs/cutcount-identity.md      the proof, the verification, the ledger
%    results/cutcount_b1/                  the banked engine source and provenance
%    results/motley-h17.md                 what the second source has closed
%    results/motley-step0.md               why the self-checks are not enough
%    polyplets/Polyplets/Defs.lean         the connectivity vocabulary
% ---------------------------------------------------------------------------

% ---------------------------------------------------------------------------
%  paper/shared/preamble.tex — packages and macros common to every manuscript
%  in this directory.  Factored out of paper/polyplets-report.tex and
%  paper/technical-report.tex, which between them had two copies of the same
%  twenty lines.
%
%  technical-report.tex does NOT \input this file and is not to be edited to do
%  so: it is jasonp's prose and is read-only to the machine (paper/README.md,
%  "Who may edit what").  The duplication there is deliberate.
%
%  Assumes \documentclass[11pt]{article} has already been issued.
% ---------------------------------------------------------------------------
\usepackage[margin=1in]{geometry}
\usepackage{amsmath,amssymb}
\usepackage{amsthm}
\usepackage{booktabs}
\usepackage{graphicx}
\usepackage{numprint}
\usepackage{microtype}
\usepackage[table]{xcolor}
\usepackage[hidelinks]{hyperref}
\usepackage{flafter}   % a float never appears before the text that introduces it
\usepackage{float}     % [H] pins tables/figures exactly where written
\usepackage[font=small,labelfont=bf,labelsep=period]{caption}
\setlength{\intextsep}{16pt plus 3pt minus 3pt}

\npthousandsep{,}
\npfourdigitsep
\newcommand{\num}[1]{\numprint{#1}}
\newcommand{\g}{\cellcolor{gray!15}}

% Theorem environments.  One counter shared across kinds, so that a reference
% like "Theorem 4" is unambiguous in a paper that also has lemmas and
% propositions.
\theoremstyle{plain}
\newtheorem{theorem}{Theorem}
\newtheorem{proposition}[theorem]{Proposition}
\newtheorem{lemma}[theorem]{Lemma}
\newtheorem{corollary}[theorem]{Corollary}
\theoremstyle{definition}
\newtheorem{definition}[theorem]{Definition}
\newtheorem{example}[theorem]{Example}
\theoremstyle{remark}
\newtheorem{remark}[theorem]{Remark}
\newtheorem{conjecture}[theorem]{Conjecture}
\newtheorem{openproblem}[theorem]{Open Problem}

% Repo-path citations.  Every one of these must resolve in the tree that ships
% with the paper; `make gate-citations` checks the markdown, and
% paper/verify_claims.py checks the manuscripts.
\newcommand{\repo}[1]{\texttt{#1}}

% OEIS links.
\newcommand{\oeis}[1]{\href{https://oeis.org/#1}{#1}}

% Growth constants, so a rename is one edit.
\newcommand{\lam}{\lambda}
\newcommand{\muH}{\mu_H}

% ---------------------------------------------------------------------------
%  paper/shared/disclosure.tex — the two authorship disclosure blocks.
%
%  docs/publication-split.md §1 fixes the rule these implement:
%
%    "Under no circumstances will there be any ambiguity about whether a paper
%     is wholly my work (none), merely my prose (some) or entirely LLM-authored
%     (some)."   — jasonp, 2026-08-07
%
%  There are exactly two blocks and they live in exactly one file, so that a
%  paper cannot quietly soften its own disclosure.  The fixed sentences take no
%  arguments.  The ONE thing a paper supplies is the per-result verification
%  ledger, because that is the part that differs honestly between papers — and
%  §1 requires it to be per result, not in aggregate.
%
%  Do not add a \Pdisclosure or \Ldisclosure variant with an optional softening
%  argument.  If a paper does not fit one of these two blocks, the paper is
%  wrong, not the block.
% ---------------------------------------------------------------------------

\newsavebox{\disclosurebox}

\newenvironment{disclosureframe}
  {\begin{center}\begin{lrbox}{\disclosurebox}\begin{minipage}{0.93\textwidth}\small}
  {\end{minipage}\end{lrbox}\fbox{\usebox{\disclosurebox}}\end{center}}

% --- P: jasonp's prose -------------------------------------------------------
% Byline: Jason Parker, alone.  Argument: the per-result ledger of what the
% machine did.
\newcommand{\Pdisclosure}[1]{%
\begin{disclosureframe}
\textbf{Authorship.}\enspace Every sentence of this paper was written by its
named author. He has read every proof it states and believes each one, and he
is responsible for the correctness of what it claims.

\smallskip
\textbf{Machine contribution.}\enspace #1
\end{disclosureframe}}

% --- L: entirely machine-written --------------------------------------------
% Argument: the per-result verification ledger — for each result, what warrants
% it (Lean kernel, exact certificate, or labelled measurement) and how much of
% it a human has checked.
\newcommand{\Ldisclosure}[1]{%
\begin{disclosureframe}
\textbf{Authorship.}\enspace This document was written end to end by a large
language model. That includes its mathematics, its prose, its tables and its
\LaTeX{}. No sentence of it was written by a human. The mathematics it reports
was likewise derived by machine.

\smallskip
\textbf{Human role.}\enspace The person who directed this work chose what to
compute, chose what to publish, and is responsible for the decision to release
this document. Except where the ledger below says otherwise, that person has
not personally checked the arguments here, and this disclosure is not to be
read as an endorsement of them.

\smallskip
\textbf{What warrants each result.}\enspace #1
\end{disclosureframe}}

% --- Draft banner ------------------------------------------------------------
% For an L-paper whose verification ledger is not yet settled.  Loud on purpose:
% an L-paper without a truthful ledger must not be mistaken for a finished one.
%
% \firstdraftbanner is the standard wording, which five papers previously
% carried character-for-character in their own preambles.  It lives here for the
% same reason the disclosure blocks do: identical text in seven places drifts,
% and a paper must not be able to soften it by editing its own copy.  The
% optional argument is for a gate specific to one paper, appended verbatim; a
% paper needing different wording (L6, compute-gated) calls \draftbanner direct.
\newcommand{\firstdraftbanner}[1][]{%
\draftbanner{This is a first draft. The verification ledger below records what
warrants each result \emph{mechanically}; the column that records what a human
has read is empty, deliberately, because nobody has read it yet. Do not
circulate until that column says something true.#1}}

\newcommand{\draftbanner}[1]{%
\begin{center}
\fcolorbox{red!60!black}{yellow!15}{%
  \begin{minipage}{0.93\textwidth}\small
  \textbf{DRAFT — NOT FOR CIRCULATION.}\enspace #1
  \end{minipage}}
\end{center}}


\newcommand{\Z}{\mathbb{Z}}
\newcommand{\Win}{\mathrm{Win}}
\newcommand{\Cfg}{\mathrm{Cfg}}

\title{\textbf{A cancellation identity for scan-order colourings}\\[4pt]
\large The rule behind a second source for lattice-animal counts}
\author{[byline: jasonp's call --- \repo{docs/publication-split.md} \S1]}
\date{August 2026}

\begin{document}
\maketitle

\firstdraftbanner[ Two gates are specific to this paper. \textbf{No
literature-priority pass has been run} (\S\ref{sec:novelty}), and the identity
is of a type with a large literature, so nothing here is claimed to be new. And
the bridge from the identity to the program that motivated it is \emph{not}
crossed here (\S\ref{sec:tier2}); that is stated as a limit, not left implicit.]

\Ldisclosure{%
\emph{Theorem~\ref{thm:main} (the identity)} --- \textbf{proved} below, for
every board and every subset, in $\Z[q]$, unconditionally. No hypothesis is
needed and none is carried; the one place a condition looked likely --- a
window-size proviso --- is discharged by Lemma~\ref{lem:reach}, which is a
theorem about the scan order rather than an assumption. Lemmas
\ref{lem:reach}--\ref{lem:choice} are elementary and self-contained. Not
formalized.
\emph{The verification of \S\ref{sec:verify}} --- exhaustive brute force over
every subset of ten boards: $223{,}296$ subsets, $331{,}935$ configurations,
zero mismatches, with five RED controls each of which fires. It reaches $16$
cells and $H \le 5$; the proof is uniform in the board, so the battery is there
to catch a mis-statement rather than to extend a range.
\emph{That the colouring model is the rule a particular program implements} ---
\textbf{checked, not proved}, and it is the single largest weak point in the
document this paper is drawn from: it is a reading of C++, argued but not
derived. A wrong reading would make the paper prove the wrong theorem, and no
internal rigour would reveal that.
\emph{That any program computes the sum} --- \textbf{not addressed}
(\S\ref{sec:tier2}).
\emph{Novelty} --- \textbf{unchecked}. See \S\ref{sec:novelty}.
\emph{Human verification:} none, as of this draft.}

\begin{abstract}
\noindent
Fix a rectangular board and scan its cells in column-major order. Colour the
cells of a subset $S$ one at a time under three local rules --- a configuration
dies if a cell sees two colours among its earlier neighbours, inherits the
colour if it sees exactly one, and otherwise may either adopt any colour
currently live in a bounded window or start a fresh one --- and weight each
surviving configuration by $\prod (q - b)$ over its fresh starts, $b$ being the
number of live colours at that moment. Then, for every board and every subset,
\[
  \sum_{\varphi} \; \prod_{v \,\in\, \mathrm{births}(\varphi)} \bigl(q - b_\varphi(v)\bigr)
  \;=\; q^{\,c(S)}
  \qquad\text{in } \Z[q],
\]
where $c(S)$ is the number of connected components of $S$. Every configuration
that fails to give each component exactly one birth cancels, exactly, with no
truncation.

The identity is proved here in five elementary lemmas. Its interest is not the
statement but what it licenses: reducing modulo $q^2$ leaves the number of
connected subsets, so a dynamic program over the window --- which never tests
connectivity --- counts connected animals. That is the mathematical content of a
second enumeration source, independent of the usual connectivity-tracking
method, and independence of method is what a count with no closed form can
otherwise never have.

What is not proved here is that any particular program computes this sum. That
bridge is stated as an open piece of work with its three parts named, rather
than assumed.
\end{abstract}

\tableofcontents

\section{Why a cancellation identity}
\label{sec:intro}

Counts of lattice animals are produced by programs, and a program's output is
believed for two reasons: it agrees with independently computed values where
they exist, and it agrees with a second program that shares no method with the
first. The second reason is the stronger one and the harder to arrange, because
the natural methods all track connectivity the same way --- by maintaining
components of a frontier and merging them --- and two programs that both do
that share their bugs.

This paper is about a rule that does not track connectivity at all. Cells are
coloured under local constraints; the weight of a colouring is a product over
the cells that started a new colour; and the sum over colourings collapses to
$q^{c(S)}$, so the coefficient of $q^1$ counts connected sets. Connectivity is
never tested. It is recovered from cancellation.

The identity below is the whole mathematical content of that rule. It is stated
for finite sets of cells and proved as combinatorics; no program appears in the
statement or in the proof.

\section{The model}
\label{sec:model}

Fix integers $H, W \ge 1$ and let $\Gamma = \{(r,c) : 0 \le r < H,\ 0 \le c < W\}$,
cells written (row, column). The \emph{scan order} is column-major:
$\mathrm{ord}(r,c) = cH + r$. Adjacency is king adjacency --- Chebyshev distance
one --- and $c(S)$ counts king-connected components of $S$, with $c(\emptyset) = 0$.

For a cell $v$ write $N^-(v)$ for its king neighbours earlier in scan order, and
let the \emph{window} $\Win(v)$ be the $H+1$ cells immediately preceding $v$.

A \emph{configuration} of $S \subseteq \Gamma$ is a map $\varphi$ from $S$ to
positive integers, built by visiting the cells of $S$ in scan order under three
rules:

\begin{description}
\item[clash.] If $\varphi(N^-(v))$ has two or more elements, the configuration is
  \emph{void} and discarded.
\item[forced.] If $\varphi(N^-(v))$ is a single label $\ell$, then
  $\varphi(v) = \ell$.
\item[free.] If $N^-(v) = \emptyset$, let
  $L_\varphi(v) = \varphi(S \cap \Win(v))$ be the \emph{live} labels and
  $b_\varphi(v) = |L_\varphi(v)|$ their number. Then either
  $\varphi(v) \in L_\varphi(v)$ (\emph{adopt}) or $\varphi(v)$ is the next unused
  label (\emph{birth}).
\end{description}

Write $\Cfg(S)$ for the non-void configurations, $B(\varphi)$ for the set of
cells at which $\varphi$ gave birth, and
\[
  w(\varphi) \;=\; \prod_{v \in B(\varphi)} \bigl(q - b_\varphi(v)\bigr) \;\in\; \Z[q],
\]
an empty product being $1$.

Three features of the free rule are load-bearing and are easy to soften by
accident.

\begin{itemize}
\item \textbf{Adopt offers every live label}, including the labels of components
  the cell does not touch. There is no adjacency test in the adopt branch.
\item \textbf{$b$ counts distinct live \emph{labels}, not components.} Two
  components sharing a colour are counted once. Because adopt lets one component
  take another's colour, the distinction is real.
\item \textbf{``Live'' means in the window.} A label whose last cell has left the
  window is gone and is not counted, which is what allows colours to be reused.
\end{itemize}

\section{The identity}
\label{sec:identity}

\begin{theorem}
\label{thm:main}
For every $H, W \ge 1$ and every $S \subseteq \Gamma$,
\[
  \sum_{\varphi \in \Cfg(S)} \ \prod_{v \in B(\varphi)} \bigl(q - b_\varphi(v)\bigr)
  \;=\; q^{\,c(S)} \qquad\text{in } \Z[q].
\]
\end{theorem}

The equality is in $\Z[q]$ itself, with no truncation --- a stronger statement
than the modulo-$q^2$ version that the application needs.

\begin{corollary}
\label{cor:count}
Summing over all $S \subseteq \Gamma$ with $|S| = n$ and reducing modulo $q^2$,
\[
  A_n(q) \;=\; \sum_{|S| = n} q^{\,c(S)}
  \;\equiv\; \#\{S : |S| = n,\ S \text{ connected}\}\cdot q \pmod{q^2},
\]
so $[q^1]A_n$ counts the connected $n$-cell subsets of the board and
$[q^0]A_n = 0$ for $n \ge 1$. Setting $q = 1$ gives $\sum_\varphi w(\varphi) = 1$
for every $S$, so the $q = 1$ evaluation counts all subsets, $\binom{HW}{n}$.
\end{corollary}

The two consequences in Corollary~\ref{cor:count} are the natural self-checks
for an implementation, and it is worth saying plainly that \emph{neither of them
consults connectivity}. A rule that has gone wrong in a way that preserves the
$q^0$ and $q=1$ evaluations passes both. Only comparison against independently
computed values catches that.

\section{Proof}
\label{sec:proof}

The proof is five lemmas. Throughout, fix $S$ and write $C_1, \dots, C_m$ for
its components, $m = c(S)$, indexed by the scan order of their minima.

\begin{lemma}[reach]
\label{lem:reach}
Every earlier king neighbour of $v$ lies in $\Win(v)$.
\end{lemma}

Under column-major order the earlier king neighbours of $(r,c)$ are at scan
distance at most $H+1$, so the window of that size sees all of them. This is a
statement about the scan order, and it is why no window-size hypothesis appears
in Theorem~\ref{thm:main}.

\begin{lemma}[survivors are constant on components]
\label{lem:const}
$\varphi \in \Cfg(S)$ if and only if $\varphi$ is constant on each component of
$S$.
\end{lemma}

\begin{lemma}[component minima are free, and nothing else is born]
\label{lem:minima}
For each $i$, the scan-minimal cell of $C_i$ is free, and every cell that gives
birth is a component minimum.
\end{lemma}

\begin{lemma}[the adopt branch a component needs is always on offer]
\label{lem:adopt}
If $v$ is free and its component already carries a label, that label is live at
$v$.
\end{lemma}

Lemma~\ref{lem:adopt} is where the window size earns its place: a component that
leaves the window and returns would have lost its label, and
Lemma~\ref{lem:reach} is exactly what rules that out.

\begin{lemma}[configurations are choice sequences]
\label{lem:choice}
The map sending $\varphi$ to its sequence of choices at the component minima is a
bijection onto the sequences in which the $i$-th minimum either adopts one of the
$b_i$ live labels or is born.
\end{lemma}

Given the lemmas, the sum factorises over the component minima: the $i$-th
contributes $b_i$ adopt terms of weight $1$ and one birth term of weight
$q - b_i$, so
\[
  \sum_{\varphi} w(\varphi)
  \;=\; \prod_{i=1}^{m} \bigl(b_i + (q - b_i)\bigr)
  \;=\; q^{\,m} \;=\; q^{\,c(S)} .
\]
The cancellation is that exact: at each component the $b_i$ ways to reuse an
existing colour are cancelled by the $-b_i$ carried in the birth weight,
leaving $q$. \qed

\section{Verification}
\label{sec:verify}

The identity is checked as stated --- by enumerating the model of
\S\ref{sec:model} directly, not by running any production code. Every subset of
each of ten boards is enumerated, weights are summed exactly in $\Z[q]$, and the
result compared with $q^{c(S)}$ using flood fill for $c(S)$.

\begin{table}[H]
\centering\small
\begin{tabular}{l r r r r}
\toprule
boards & subsets & configurations & mismatches & wall \\
\midrule
$1{\times}6$, $1{\times}10$, $2{\times}5$, $2{\times}7$, $2{\times}8$,
 & $223{,}296$ & $331{,}935$ & $0$ & $5.4$\,s \\
$3{\times}4$, $3{\times}5$, $4{\times}3$, $4{\times}4$, $5{\times}3$ & & & & \\
\bottomrule
\end{tabular}
\caption{Exhaustive over subsets, not sampled. Lemma~\ref{lem:adopt} is
additionally audited at $53{,}552$ instances, with zero failures.}
\label{tab:verify}
\end{table}

Five RED controls run alongside, and all five fire:

\begin{table}[H]
\centering\small
\begin{tabular}{l l}
\toprule
control & what it corrupts \\
\midrule
RED-1 & clash-voiding dropped \\
RED-2 & $b$ counts all labels, not live ones \\
RED-3 & a free cell may not adopt a live label \\
RED-4 & window of $H$ cells, so some earlier neighbours are unseen \\
RED-5 & liveness window shortened to $H-2$ \\
\bottomrule
\end{tabular}
\caption{Each control breaks one ingredient the proof uses, and each is required
to produce a mismatch. RED-4 and RED-5 are the ones that put
Lemma~\ref{lem:reach} under test.}
\label{tab:red}
\end{table}

Two limits of this battery, stated rather than left to be discovered. It reaches
$16$ cells and $H \le 5$, so a phenomenon first appearing at $H \ge 6$ would not
be seen; the proof is uniform in $H$ and $W$, which is why a bounded check is the
right shape of test, but it is not an unbounded one. And the controls test the
ingredients identified as load-bearing --- an ingredient nobody thought of has no
control.

\section{What is not proved here}
\label{sec:tier2}

The identity is about finite sets of cells. That some \emph{program} computes the
sum is a separate claim, and this paper does not make it. The argument has three
parts, all mechanical and none carried out here:

\begin{itemize}
\item \textbf{Locality} --- the window sees all earlier neighbours. This half is
  proved: it is Lemma~\ref{lem:reach}.
\item \textbf{Sufficient statistic} --- the rules consult the past only through
  equality of labels within the window and the number of distinct live labels, so
  the colour-coincidence partition of the window determines the remaining weight,
  and a canonical form of that partition is a sufficient state.
\item \textbf{Colour reuse} --- an implementation's fresh label is fresh only
  within the window, whereas \S\ref{sec:model}'s births are globally fresh. The
  two agree because an expired label is never consulted again.
\end{itemize}

Beyond that sits the bridge from a formal model to a compiled binary, which is
not addressed at all. And one further layer is disjoint from everything here: the
accounting that turns per-board counts into counts by bounding-box height. The
identity is per-board and says nothing about it.

\section{Novelty}
\label{sec:novelty}

\textbf{Unchecked.} No literature search has been run on this statement, and
unlike most of this series the base rate is unfavourable: counting connected
subgraphs by a colouring with signed weights, so that disconnected configurations
cancel, is a well-populated idea, and the transfer-matrix literature on lattice
animals is large. The honest expectation is that something close to
Theorem~\ref{thm:main} exists. Nothing here is claimed to be new; what is claimed
is that this particular rule, with this window and this liveness convention, is
proved to compute $q^{c(S)}$.

\section{Open problems}

\begin{openproblem}[formalize the identity]
Lemmas~\ref{lem:reach}--\ref{lem:choice} are elementary, finite in statement and
uniform in the board. The repository's Lean development already carries the
connectivity vocabulary and a comparable chain, so this is a scoped target rather
than a research problem.
\end{openproblem}

\begin{openproblem}[cross the bridge of \S\ref{sec:tier2}]
The three parts are stated. Doing them turns ``a program was checked against
known values'' into ``a program provably computes a proved identity'', which is
the difference the second source exists to make.
\end{openproblem}

\begin{openproblem}[other adjacencies and other scan orders]
Nothing in the proof is specific to king adjacency beyond
Lemma~\ref{lem:reach}, which bounds the scan distance to earlier neighbours. For
which adjacency and scan-order pairs does a window of bounded size suffice, and
what is the minimal window?
\end{openproblem}

\section{Reproduction}
\label{sec:repro}

\begin{table}[H]
\centering\small
\begin{tabular}{l l}
\toprule
claim & check \\
\midrule
the identity, on every subset of ten boards & \repo{experiments/birthright\_identity\_check.py} \\
the five RED controls & same script \\
Lemma~\ref{lem:adopt}, audited & same script \\
the proof, in full & \repo{docs/proofs/cutcount-identity.md} \\
the banked rule the model reads & \repo{results/cutcount\_b1/} \\
what the second source has closed & \repo{results/motley-h17.md} \\
why the self-checks are not sufficient & \repo{results/motley-step0.md} \\
\bottomrule
\end{tabular}
\caption{The verification script is pure Python and takes seconds; it enumerates
the model of \S\ref{sec:model} and touches no production engine.}
\label{tab:repro}
\end{table}

\bibliographystyle{plain}
\bibliography{shared/refs}

\end{document}
